\part{统计与概率}
这一部分分为四章
\chapter{统计（一）}


\section{数据的收集}
\subsection{总体、样本与抽样}
\clearpage
\subsubsection{如何统计？}
\clearpage
\subsubsection{抽样方法}
\clearpage


\section{数据的可视化}
\subsection{条形统计图、折线统计图与扇形统计图}
\clearpage
\subsection{直方图与茎叶图}
\clearpage


\section{样本的数字特征}
\subsection{平均数}
\clearpage
\subsection{方差}
\clearpage
\subsection{\texorpdfstring{$p$}a分位数}
\clearpage
\subsection{用样本估计总体}
\clearpage
\subsection{其他重要的数字特征}
\clearpage


\section{计数原理}
\subsection{加法原理与乘法原理}
\begin{theorem}[分类加法计数原理]
    完成一件事，如果有 $n$ 类办法，且：第一类办法中有 $m_1$ 种不同的方法，第二类办法中有 $m_2$ 种不同的方法……第 $n$ 类办法中有 $m_n$ 种不同的方法，那么完成这件事共有
    \[ N = m_1 + m_2 + \cdots + m_n \]
    种不同的方法。
\end{theorem}
\begin{theorem}
    完成一件事，如果需要分成 $n$ 个步骤，且：做第一步有 $m_1$ 种不同的方法，做第二步有 $m_2$ 种不同的方法……做第 $n$ 步有 $m_n$ 种不同的方法，那么完成这件事共有
\[ N = m_1 \times m_2 \times \cdots \times m_n \]
\end{theorem}
\clearpage
\subsection{排列数与组合数}
\begin{definition}
    一般地，从 \( n \) 个不同对象中，任取 \( m \) (\( m \leq n \)) 个对象，按照一定的顺序排成一列，
    称为从 \( n \) 个不同对象中取出 \( m \) 个对象的一个\textcolor{blue}{排列}。特别地，\( m = n \) 时的排列（即取出所有对象的排列）称为\textcolor{blue}{全排列}。
    从 \( n \) 个不同对象中取出 \( m \) 个对象的所有排列的个数，称为从 \( n \) 个不同对象中取出 \( m \) 个对象的排列数，用符号 \( \mathrm{A}_n^m \) 表示。
\end{definition}

\begin{theorem}
    \[\mathrm{A}_n^m=\underbrace{ n(n-1)\cdots[n-(m-1)]}_{n\text{个连续整数}}=\frac{n!}{(n-m)!}\]
\end{theorem}

\begin{example}
    求证：\( A_n^m + mA_n^{m-1} = A_{n+1}^m \)。

\end{example}
\begin{proof}
    \begin{way1}
        由排列数公式可知
\begin{align*}
A_n^m + mA_n^{m-1} &= \frac{n!}{(n-m)!} + m \frac{n!}{[n-(m-1)]!} \\
&= \frac{n!}{(n-m)!} \times \left[1 + \frac{m}{n-(m-1)}\right] \\
&= \frac{n!}{(n-m)!} \times \frac{n+1}{n-(m-1)} \\
&= \frac{(n+1)!}{[(n+1)-m]!} = A_{n+1}^m.
\end{align*}
    \end{way1}

    \begin{way2}
        第一类：不包括甲的排列。这种情况下，我们实际上是从剩下的 \( n \) 个对象中取出 \( m \) 个进行排列，所以排列数是 \( A_n^m \)。

第二类：包括甲的排列。这里，甲必须在排列中的某个位置。既然甲可以占据 \( m \) 个位置中的任何一个，那么对于每一个位置，甲固定在这个位置后，剩下的 \( m - 1 \) 个位置需要从剩下的 \( n \) 个对象中取出并排列。因此，排列数是 \( m \times A_n^{m-1} \)。

将这两类排列数相加，就得到了总的排列数：
\[ A_n^m + mA_n^{m-1} \]
    \end{way2}
\qed
\end{proof}

\begin{definition}
    从 \( n \) 个不同对象中取出 \( m \) (\( m \leq n \)) 个对象并成一组，称为从 \( n \) 个不同对象中取出 \( m \) 个对象的一个组合。

从 \( n \) 个不同对象中取出 \( m \) 个对象的所有组合的个数，称为从 \( n \) 个不同对象中取出 \( m \) 个对象的组合数，用符号 \( C_n^m \) 表示。
\end{definition}
\begin{theorem}
    \[
C_n^m = \frac{A_n^m}{A_m^m} = \frac{n(n-1)\cdots[n-(m-1)]}{m \times (m-1) \times \cdots \times 2 \times 1} = \frac{n!}{(n-m)!m!}.
\]
\end{theorem}
\begin{proof}
    考虑到从 \( n \) 个不同对象中取出 \( m \) 个做排列，可以分成两个步骤来完成：第一步，从 \( n \) 个不同对象中取出 \( m \) 个，有 \( C_n^m \) 种选法；第二步，将选出的 \( m \) 个对象做全排列，有 \( A_m^m \) 种排法。由分步乘法计数原理有 \( A_n^m = C_n^m A_m^m \)，所以
\[
C_n^m = \frac{A_n^m}{A_m^m} = \frac{n(n-1)\cdots[n-(m-1)]}{m \times (m-1) \times \cdots \times 2 \times 1} = \frac{n!}{(n-m)!m!}.
\]
\end{proof}
\begin{example}
    证明下列组合数公式：
    \[ C_n^{m+1} + C_n^m = C_{n+1}^{m+1}\]
    \[\mathrm{C}_n^k=\mathrm{C}_n^{n-k}\]

\end{example}

\begin{example}证明下列组合恒等式：

    \begin{enumerate}
        \item \[C_n^m = \frac{n-m+1}{m} C_n^{m-1}\]
        \item \[\sum_{i=r}^{n} C_i^r = C_{n+1}^{r+1}\]
        \item \[\sum_{i=0}^{r} C_m^{r-i} C_n^i = C_{m+n}^r\]
    \end{enumerate}
\end{example}

\begin{proof}
    \begin{enumerate}
        \item 现有 \( n \) 个球，求从中挑出 \( m \) 个且把其中一个涂红有几种不同方案？

        算一次：直接 \( n \) 取 \( m \)，再在 \( m \) 个中挑一个出来染红，共 \( m C_n^m \) 种
        
        算两次：先 \( n \) 取 \( m-1 \)，再在剩下 \( n-m+1 \) 个挑一个出来并染红，共 \( (n-m+1) C_n^{m-1} \)
        \item 现有 \( n+1 \) 个球，求从中挑出 \( r+1 \) 个有几种不同方案？

        算一次：\( C_{n+1}^{r+1} \)
        
        算两次：将球标上 1 到 \( n+1 \) 号，
        
        包含 1 号球的选法有 \( C_n^r \)，不包含 1 号球包含 2 号球的选法有 \( C_{n-1}^r \)，不包含 1、2 号球包含 3 号球的选法有 \( C_{n-2}^r \)，依此类推，总共 \( \sum_{i=r}^{n} C_i^r \)
        \item 现有 \( m+n \) 个球，求从中挑出 \( r \) 个有几种不同方案？

        算一次：\( C_{m+n}^r \)
        
        算两次：任意分成两批，一批 \( m \) 个球，一批 \( n \) 个球，从第一批选 \( r-i \) 个，第二批选 \( i \) 个共 \( C_m^{r-i} C_n^i \) 种选法，总共有 \( \sum_{i=0}^{r} C_m^{r-i} C_n^i \)
    \end{enumerate}
\end{proof}





\clearpage
\subsection{排列组合的综合应用}













顺序固定
\begin{exercise}
    某货场有两堆集装箱，一堆 2 个，一堆 3 个，现需要全部装运，每次只能从其中一堆取最上面的一个集装箱，则在装运的过程中不同取法的种数.
\end{exercise}



瑞士数学家欧拉按一般情况给出了一个递推公式：用 $A$、$B$、$C$、$\cdots$ 表示写着 $n$ 位友人名字的信封，$a$、$b$、$c$、$\cdots$ 表示 $n$ 份相应的写好的信纸。把错装的总数为记作 $f(n)$。假设把 $a$ 错装进 $B$ 里了，包含着这个错误的一切错装法分两类：
\begin{enumerate}
    \item $b$ 装入 $A$ 里，这时每种错装的其余部分都与 $A$、$B$、$a$、$b$ 无关，应有 $f(n-2)$ 种错装法。
    \item $b$ 装入 $A$、$B$ 之外的一个信封，这时的装信工作实际是把（除 $a$ 之外的）$n-1$ 份信纸 $b$、$c$、$\cdots$ 装入（除 $B$ 以外的）$n-1$ 个信封 $A$、$C$、$\cdots$，显然这时装错的方法有 $f(n-1)$ 种。
\end{enumerate}
总之在 $a$ 装入 $B$ 的错误之下，共有错装法 $f(n-2) + f(n-1)$ 种。
$a$ 装入 $C$，装入 $D$、$\cdots$ 的 $n-2$ 种错误之下，
同样都有 $f(n-2) + f(n-1)$ 种错装法。

因此，得到一个递推公式：$f(n) = (n-1)\{f(n-2) + f(n-1)\}$，分别代入 $n = 2, 3, 4$ 等，可推得结果。

也可用迭代法推导出一般公式：$f(n) = n!\left(1 - \frac{1}{1!} + \frac{1}{2!} - \frac{1}{3!} + \cdots + (-1)^n \frac{1}{n!}\right)$。











\clearpage

\review

\hw
\begin{homework}
    三个比赛项目，六人报名参加。

（1）每人参加一项有多少种不同的方法？

（2）每项一人，且每人至多参加一项，有多少种不同的方法？

（3）每项一人，每人参加的项数不限，有多少种不同的方法？
\end{homework}
\begin{homework}
 用 0，1，2，3，4，5 六个数字

（1）可以组成多少个三位数？

（2）可以组成多少个无重复数字的三位数？

（3）可以组成多少个无重复数字的三位偶数？

（4）可以组成多少个无重复数字的能被 $5$ 整除的三位数？

（5）可以组成多少个无重复数字的大于 $300 $的三位偶数？
\end{homework}
\begin{homework}
    一个口袋内装有大小相同且编有不同号码的$ 5 $个白球和$ 4 $个黑球。


（1）从口袋内取出 $3 $个球，共有多少种取法？

（2）从口袋内取出 $3$ 个球，使其中恰有 $1 $个黑球，共有多少种取法？

（3）从口袋内取出 $3 $个球，使其中至少有$ 1 $个黑球，共有多少种取法？
\end{homework}
\begin{homework}
    从 $7$ 名男生 $5 $名女生中，选出 5 人，分别求符合下列条件的选法种数有多少种？

（1）A、B 必须当选；

（2）A、B 都不当选；

（3）A、B 不全当选；

（4）至少有 $2$ 名女生当选；

（5）选出$ 5 $名同学，让他们分别担任体育委员、文娱委员等 $5$ 种不同工作，但体育委员由男生担任，文娱委员由女生担任。
\end{homework}
\begin{homework}
    有$ 5$ 个男生和 $3 $个女生，从中选取 $5 $人担任$ 5 $门不同学科的科代表，求分别符合下列条件的选法数：

（1）有女生但人数必须少于男生。

（2）某女生一定要担任语文科代表。

（3）某男生必须包括在内，但不担任数学科代表。

（4）某女生一定要担任语文科代表，某男生必须担任科代表，但不担任数学科代表。
\end{homework}
\begin{homework}
现在有\(12\)个相同的小球：

（1）放入编号为 1，2，3，4 的盒子中，问每个盒子中至少有一个小球的不同放法有多少种？

（2）放入编号为 1，2，3，4 的盒子中，要求每个盒子中的小球数不小于其编号数，问不同的放法有多少种？

（3）放入编号为 1，2，3，4 的盒子中，每盒可空，问不同的放法有多少种？
\end{homework}
\begin{homework}
    在高三（1）班进行的演讲比赛中，共有 5 位选手参加，其中 3 位女生，2 位男生。如果 2 位男生不能连续出场，且女生甲不能排在第一个，那么出场顺序的排法种数
\end{homework}
\begin{homework}
    8 名运动员参加男子 100 米的决赛。已知运动场有从内到外编号依次为 1，2，3，4，5，6，7，8 的八条跑道，若指定的 3 名运动员所在的跑道编号必须是三个连续数字（如：4，5，6），则参加比赛的这 8 名运动员安排跑道的方式
\end{homework}
\begin{homework}
    在航天员进行一项太空实验中，要先后实施 6 个程序，其中程序 A 只能出现在第一或最后一步，程序 B 和 C 在实施时必须相邻，问实验顺序的编排方法共有
\end{homework}
\begin{homework}
    若有 5 本不同的书，分给三位同学，每人至少一本，则不同的分法数是（ ）
\end{homework}
\begin{homework}
    某学校需要把 6 名实习老师安排到 A，B，C 三个班级去听课，每个班级安排 2 名老师，已知甲不能安排到 A 班，乙和丙不能安排到同一班级，则安排方案的种数有（ ）

\end{homework}
\begin{homework}
    在高校自主招生中，某学校获得 5 个推荐名额，其中清华大学 2 名，北京大学 2 名，复旦大学 1 名。并且北京大学和清华大学都要求必须有男生参加。学校通过选拔定下 3 男 2 女共 5 个推荐对象，则不同的推荐方法共有（ ）
\end{homework}
\begin{homework}
    从 5 名男生医生、2 名女医生中选 3 名医生组成一个医疗小分队，要求其中至少有 1 名女医生，则不同的组队方案共有（ ）
\end{homework}
\begin{homework}
    有甲、乙、丙三项任务，甲需 2 人承担，乙、丙各需 1 人承担，从 10 人中选派 4 人承担这三项任务的不同选法有 种。
\end{homework}
\begin{homework}
    从\(6\)名运动员中选出\(4\)人参加 \(4\times 100\)米接力赛，如果甲不跑第一棒，乙不跑第四棒，共有多少种不同的参赛方案？
\end{homework}
\begin{homework}
    
\end{homework}
\begin{homework}
    
\end{homework}
\begin{homework}
    
\end{homework}




\clearpage


\section{二项式定理}
\subsection{二项式定理}
\clearpage
\subsection{多项式定理*}
\clearpage


\chapter{概率（一）}


\section{随机事件与样本空间}
\subsection{样本空间}
\clearpage
\subsection{事件的关系和运算}
\clearpage
\subsection{概率的基本性质}
公理化定义 加法公式
\clearpage
\subsection{频率与概率的关系}
\clearpage


\section{古典概型与几何概型}
\subsection{古典概型}
\clearpage
\subsection{几何概型*}
\clearpage

\section{条件概率}
\subsection{条件概率}
\clearpage
\subsection{事件的独立性}
\clearpage
\subsection{全概率公式与Bayes公式}
\subsubsection{全概率公式}
\clearpage
\subsubsection{Bayes公式}
\clearpage


\chapter{概率（二）}
\section{随机变量}
\subsection{随机变量的定义}
\clearpage
\subsection{随机变量的分布列与分布函数}
\clearpage


\section{重要的几种分布}
\subsection{两点分布与二项分布}
\clearpage
\subsection{超几何分布}
\clearpage
\subsection{正态分布}
\clearpage
\subsection{Bernoulli试验诱导的其他常见分布*}
\clearpage


\section{随机变量的数字特征}
\subsection{离散随机变量的期望与方差}
期望的重要性质，期望与方差的计算
\clearpage
\subsection{连续随机变量的期望与方差**}
\clearpage
\subsection{协方差与相关系数**}
\clearpage
\subsection{大数定律与中心极限定理简介**}
\clearpage


\section{Markov链*}
\clearpage


\chapter{统计（二）}
\section{数理统计中常见的分布**}
\subsection{三大分布}
\clearpage
\subsection{次序统计量}
\clearpage



\section{回归分析}
\subsection{最小二乘法与线性回归}
\clearpage
\subsection{非线性回归}
\clearpage



\section{显著性检验**}
\subsection{Fisher显著性检验}
\clearpage
\subsection{单样本正态总体参数的显著性检验}
\clearpage


\section{独立性检验}
\subsection{\texorpdfstring{$\chi^2$}a检验}
\clearpage
\subsection{其他独立性检验方法简介**}
\clearpage

